\chapter{Mecanismos de reação dos complexos}\label{ch:b3:complex-mechanisms}

Dissolva um sal de crômio(III) em água marcada com oxigênio-18 e passam-se dias
até que a água marcada substitua as moléculas de água ligadas ao
metal; com um sal de cobre(II), a troca se completa no tempo da
mistura. A reação é a mesma, uma molécula de água deixando um íon metálico e
outra tomando seu lugar, mas sua velocidade varia enormemente de um íon para o
outro, e o número de elétrons $d$ prevê boa parte dessa variação. A maneira como
os ligantes são substituídos decide como os fármacos anticâncer de platina são preparados e como
agem; a maneira como os elétrons passam entre íons metálicos decide a velocidade
da respiração e de toda bateria. Este capítulo trata das duas famílias de
reações dos complexos: a substituição e a transferência de elétrons, esta última com
a teoria de Rudolph Marcus.

\begin{recall}
O volume do 2.º ano descreveu os complexos, a denticidade dos ligantes, os ligantes
em ponte, os isômeros $fac$ e $mer$, a troca de ligantes como etapa elementar e as
energias de estabilização do campo cristalino com as configurações de spin alto e baixo;
o volume do 1.º ano, os mecanismos com as aproximações do estado estacionário e da etapa
determinante. O \cref{ch:b3:rate-theories} deu a \omterm{thm:b3:rate-theories:eyring}{equação de Eyring} e o
significado de $\Delta^{\ddagger}S^\circ$, o \cref{ch:b3:electrode-kinetics}, a transferência de elétrons em um
eletrodo, e o \cref{ch:b3:complex-spectra-magnetism}, os \omterm{def:b3:complex-spectra-magnetism:ligand-field-term}{termos do campo ligante} e o
efeito Jahn--Teller.
\end{recall}

\section{Labilidade e inércia}

\begin{definition}[Complexos lábeis e inertes]\label{def:b3:complex-mechanisms:labile}
Um complexo é \emph{lábil}\index{complexo labil@complexo lábil} se seus ligantes são substituídos
rapidamente, e um \emph{complexo cineticamente inerte}\index{complexo cineticamente inerte}
se são substituídos lentamente; seguindo Taube, um complexo é dito lábil quando
suas reações com um ligante a \qty{0.1}{mol/L} e \qty{25}{\celsius} terminam em cerca de um
minuto. A labilidade é uma propriedade cinética, sem relação com a estabilidade
termodinâmica do complexo.
\end{definition}

A troca de água, $\mathrm{[M(H_2O)_6]^{n+} + H_2O^{\ast} \to [M(H_2O)_5(H_2O^{\ast})]^{n+} + H_2O}$, é a substituição mais
simples e serve de escala. Os íons dos grupos principais trocam mais depressa quando são
maiores e menos carregados. Entre os íons de metais de transição, os de configuração $d^3$
(\ce{Cr^{3+}}) ou $d^6$ de spin baixo (\ce{Co^{3+}}, \ce{Rh^{3+}}, \ce{Ir^{3+}}) são inertes, assim como
a maioria dos metais da segunda e da terceira séries, ao passo que $d^9$ (\ce{Cu^{2+}}) e $d^4$ de spin alto (\ce{Cr^{2+}}),
distorcidos pelo efeito Jahn--Teller, são extremamente \omterm{def:b3:complex-mechanisms:labile}{lábeis}.

\begin{proposition}[Energia de ativação do campo ligante]\label{prop:b3:complex-mechanisms:lfae}
Se uma substituição passa por um estado de transição pentacoordenado piramidal de
base quadrada, a perda de estabilização do campo ligante ao chegar ali (a energia de
ativação do campo ligante) é, no modelo mais simples, só $\sigma$, máxima para $d^6$ de spin baixo
($4\,Dq$) e depois para $d^3$ e $d^8$ ($2\,Dq$), nula para $d^0$, $d^5$ de spin alto e $d^{10}$, e negativa
para $d^4$ de spin alto e $d^9$.
\end{proposition}

\begin{proof}
Modele cada ligante como elevando um orbital $d$ de $e_\sigma$ vezes o quadrado da amplitude
angular do orbital em sua direção, normalizada a 1 no eixo do lobo: um ligante
axial dá a $d_{z^2}$ 1, um equatorial dá a $d_{z^2}$ $\frac14$ e a $d_{x^2-y^2}$ $\frac34$; os orbitais $t_{2g}$ apontam
entre os ligantes e não recebem nada. O octaedro dá a $d_{z^2}$ e a $d_{x^2-y^2}$ $3e_\sigma$ cada
($\Delta_{\mathrm o} = 3e_\sigma = 10Dq$); a pirâmide de base quadrada, sem um ligante axial, dá a $d_{z^2}$ $2e_\sigma$ e a
$d_{x^2-y^2}$ $3e_\sigma$. A estabilização de $n$ elétrons é sua energia menos $n/5$ vezes a soma dos
cinco níveis (o baricentro). Para $d^6$ de spin baixo ($t_{2g}^6$): $0 - \frac65 \times 6e_\sigma$ no
octaedro, $0 - \frac65 \times 5e_\sigma$ na pirâmide: uma perda de $1.2e_\sigma = 4Dq$. Para $d^3$: $\frac35e_\sigma =
2Dq$; para $d^8$ ($t_{2g}^6e_g^2$): $(5 - \frac85 \times 5) - (6 - \frac85 \times 6) = 0.6e_\sigma$. Para $d^9$: $(7 - 9) - (9 -
10.8) = -0.2e_\sigma$.
\end{proof}

\begin{omfigure}
\begin{tikzpicture}
\begin{axis}[width=0.72\linewidth, height=5.4cm, xmin=-0.6, xmax=10.6, ymin=-1.2, ymax=4.6,
  axis x line=bottom, axis y line=left, ycomb, xtick={0,1,2,3,4,5,6,7,8,9,10},
  xlabel={$n$ ($d^n$)}, ylabel={energia de ativação / $Dq$}, tick label style={font=\small},
  label style={font=\small}]
\draw[black!40] (axis cs:-0.6,0) -- (axis cs:10.6,0);
\addplot[omDef, line width=4pt, xshift=-2.5pt] table[x=n, y=hs] {figdata/out/bachelor-3-ligand-field-activation.dat};
\addplot[omThm, line width=4pt, xshift=2.5pt, unbounded coords=discard] table[x=n, y=ls] {figdata/out/bachelor-3-ligand-field-activation.dat};
\end{axis}
\end{tikzpicture}

{\small Energias de ativação do campo ligante para um caminho dissociativo por uma
pirâmide de base quadrada, no modelo só $\sigma$ da proposição (azul: spin alto;
vermelho: spin baixo). Os maiores valores,
para $d^6$ de spin baixo, $d^3$ e $d^8$, correspondem aos íons inertes ou lentos (\ce{Co^{3+}}, \ce{Cr^{3+}}, \ce{Ni^{2+}});
os negativos, para $d^4$ e $d^9$, aos íons de Jahn--Teller muito \omterm{def:b3:complex-mechanisms:labile}{lábeis}.}
\end{omfigure}

\begin{definition}[Volume de ativação]\label{def:b3:complex-mechanisms:activation-volume}
O \emph{volume de ativação}\index{volume de ativacao@volume de ativação} $\Delta V^{\ddagger}$ de uma etapa
elementar é a diferença entre o volume parcial molar do \omterm{def:b3:rate-theories:activated-complex}{complexo
ativado} e os dos reagentes.
\end{definition}

\begin{proposition}[Dependência de uma constante de velocidade com a pressão]\label{prop:b3:complex-mechanisms:pressure}
A temperatura constante, $\displaystyle\Bigl(\frac{\partial\ln k}{\partial p}\Bigr)_T = -\frac{\Delta V^{\ddagger}}{RT}$.
\end{proposition}

\begin{proof}
Pela \omterm{thm:b3:rate-theories:eyring}{equação de Eyring}, $\ln k = \text{const} - \Delta^{\ddagger}G/RT$, e $(\partial\Delta^{\ddagger}G/\partial p)_T = \Delta V^{\ddagger}$, como para
toda energia de Gibbs.
\end{proof}

Uma etapa na qual um ligante sai tem $\Delta V^{\ddagger} > 0$; uma na qual um ligante entra tem $\Delta V^{\ddagger}
< 0$. Medido a algumas centenas de megapascals, o \omterm{def:b3:complex-mechanisms:activation-volume}{volume de ativação} é o diagnóstico mais
direto de um mecanismo de substituição, menos obscurecido pela solvatação do que
a \omterm{def:b3:rate-theories:activation}{entropia de ativação}.

\section{Mecanismos de substituição}

\begin{definition}[Mecanismos dissociativo, associativo e de intercâmbio]\label{def:b3:complex-mechanisms:mechanisms}
Em um \emph{mecanismo dissociativo}\index{mecanismo dissociativo} (D), o ligante
de saída parte primeiro, dando um intermediário de número de coordenação menor; em um
\emph{mecanismo associativo}\index{mecanismo associativo} (A), o ligante de entrada
se liga primeiro, dando um intermediário de número de coordenação maior; em um
\emph{mecanismo de intercâmbio}\index{mecanismo de intercambio@mecanismo de intercâmbio} (I), os dois ligantes
se trocam em uma única etapa, com um caráter de ruptura de ligação ($I_d$) ou de formação de ligação ($I_a$).
\end{definition}

\begin{proposition}[Leis de velocidade dos mecanismos D e A]\label{prop:b3:complex-mechanisms:d-a-laws}
Para $\mathrm{ML_5X + Y \to ML_5Y + X}$: um mecanismo D ($\mathrm{ML_5X} \rightleftharpoons \mathrm{ML_5} + \mathrm X$, $k_1$, $k_{-1}$; $\mathrm{ML_5} + \mathrm Y
\to \mathrm{ML_5Y}$, $k_2$) dá
\[
  v = \frac{k_1k_2[\mathrm{ML_5X}][\mathrm Y]}{k_{-1}[\mathrm X] + k_2[\mathrm Y]},
\]
de primeira ordem no complexo e independente de $[\mathrm Y]$ para $[\mathrm Y]$ alto; um mecanismo A
por um intermediário heptacoordenado em estado estacionário dá $v =
k[\mathrm{ML_5X}][\mathrm Y]$, de primeira ordem em cada um.
\end{proposition}

\begin{proof}
D: estado estacionário para $\mathrm{ML_5}$, $k_1[\mathrm{ML_5X}] = (k_{-1}[\mathrm X] + k_2[\mathrm Y])[\mathrm{ML_5}]$, e $v = k_2[\mathrm{ML_5}][\mathrm Y]$.
Para $[\mathrm Y]$ alto, $v \to k_1[\mathrm{ML_5X}]$. A: $\mathrm{ML_5X} + \mathrm Y \rightleftharpoons \mathrm{ML_5XY}$ ($k_a$, $k_{-a}$), $\mathrm{ML_5XY} \to
\mathrm{ML_5Y} + \mathrm X$ ($k_b$); o estado estacionário dá $v = \frac{k_ak_b}{k_{-a} + k_b}[\mathrm{ML_5X}][\mathrm Y]$.
\end{proof}

Na água, o ligante de entrada costuma estar presente em concentração muito menor
que o solvente, e a verdadeira primeira etapa é o encontro do complexo com o
ligante.

\begin{definition}[Mecanismo de Eigen--Wilkins]\label{def:b3:complex-mechanisms:eigen-wilkins}
O \emph{mecanismo de Eigen--Wilkins}\index{mecanismo de Eigen--Wilkins} de
substituição de um aquacomplexo é um pré-equilíbrio rápido que forma um
\emph{complexo de esfera externa}\index{complexo de esfera externa}, no qual o ligante
de entrada fica junto ao complexo, fora de sua esfera de coordenação, seguido do
intercâmbio, determinante da velocidade, de uma molécula de água por esse ligante.
\end{definition}

\begin{theorem}[Lei de velocidade de Eigen--Wilkins]\label{thm:b3:complex-mechanisms:eigen-wilkins}
Com a constante de equilíbrio de esfera externa $K_{\mathrm{os}}$ e a constante de velocidade de intercâmbio
$k_i$, a constante de velocidade de pseudoprimeira ordem observada com $[\mathrm Y]$ em excesso é
\[
  k_{\mathrm{obs}} = \frac{K_{\mathrm{os}}k_i[\mathrm Y]}{1 + K_{\mathrm{os}}[\mathrm Y]} .
\]
\end{theorem}

\begin{proof}
O complexo se distribui entre as formas livre e de esfera externa, $[\mathrm{M\cdot Y}] = K_{\mathrm{os}}[\mathrm M][\mathrm Y]$,
com total $[\mathrm M]_{\mathrm{tot}} = [\mathrm M](1 + K_{\mathrm{os}}[\mathrm Y])$; a velocidade é $k_i[\mathrm{M\cdot Y}] = k_iK_{\mathrm{os}}[\mathrm Y][\mathrm M]_{\mathrm{tot}}/(1 + K_{\mathrm{os}}[\mathrm Y])$.
\end{proof}

Para $[\mathrm Y]$ baixo, a reação é de segunda ordem com $k = K_{\mathrm{os}}k_i$; $K_{\mathrm{os}}$, estimada a partir das
cargas e da distância de máxima aproximação, é quase a mesma para todos os
ligantes de mesma carga, de modo que $k_i$, a etapa de troca de água, controla a
velocidade: a substituição em um aquaíon é quase independente do ligante
de entrada, a assinatura de um intercâmbio dissociativo.

\begin{definition}[Mecanismo da base conjugada]\label{def:b3:complex-mechanisms:conjugate-base}
O \emph{mecanismo da base conjugada}\index{mecanismo da base conjugada} da hidrólise
básica de um aminocomplexo, $\mathrm{[Co(NH_3)_5X]^{2+} + OH^- \to [Co(NH_3)_5(OH)]^{2+} + X^-}$, é
a desprotonação de um ligante amin em um pré-equilíbrio rápido, seguida da
perda dissociativa de $\mathrm X^-$ pela base conjugada amido, cujo ligante $\mathrm{NH_2^-}$
labiliza o complexo.
\end{definition}

\begin{proposition}[Lei de velocidade da hidrólise básica]\label{prop:b3:complex-mechanisms:conjugate-base}
Com a constante de desprotonação $K$ e a constante de velocidade de dissociação $k$ da
base conjugada, $v = \dfrac{kK[\text{complexo}]_{\mathrm{tot}}[\ce{OH-}]}{1 + K[\ce{OH-}]}$, de segunda ordem para $[\ce{OH-}]$ baixo.
\end{proposition}

\begin{proof}
Como para a lei de Eigen--Wilkins: o complexo se distribui entre suas formas ácida e
de base conjugada na razão $1 : K[\ce{OH-}]$, e só a segunda reage.
\end{proof}

A lei de velocidade sozinha não distingue esse mecanismo de um ataque direto do
\ce{OH-}; as evidências são que a hidrólise básica exige um próton N--H (os complexos sem
ele não a apresentam) e que os prótons dos ligantes amin trocam com o solvente muito
mais depressa que a hidrólise.

\begin{method}[Diagnosticar um mecanismo de substituição]\label{met:b3:complex-mechanisms:diagnose}
\begin{enumerate}
  \item Dependência do grupo de entrada: velocidade independente de Y (ou que satura),
        dissociativo; velocidade proporcional a $[\mathrm Y]$ e sensível à natureza de Y,
        associativo.
  \item \omterm{def:b3:complex-mechanisms:activation-volume}{Volume de ativação}: positivo, dissociativo; negativo, associativo.
  \item \omterm{def:b3:rate-theories:activation}{Entropia de ativação}: positiva para D, negativa para A (com cautela, pois a
        solvatação contribui).
  \item Dependência do grupo de saída: uma velocidade que acompanha a força da ligação M--X
        indica ruptura de ligação no estado de transição.
\end{enumerate}
\end{method}

\begin{omfigure}
\begin{tikzpicture}
\begin{axis}[name=pd, omprofile, width=0.46\linewidth, height=4.6cm, xmin=0, xmax=10, ymin=0, ymax=10,
  xlabel={coordenada de reação}, ylabel={energia}, title={\small D: $\mathrm{ML_5X} \to \mathrm{ML_5} + \mathrm X \to \mathrm{ML_5Y}$}]
\addplot[omDef, very thick, smooth] coordinates {(0,2) (1.5,2.1) (3,7.5) (4.2,5.6) (5.0,5.4) (5.8,5.6) (7,7.0) (8.5,1.6) (10,1.5)};
\node[font=\scriptsize, anchor=north, align=center] at (axis cs:5,5.2) {ML$_5$ (NC 5)};
\end{axis}
\begin{axis}[at={(pd.east)}, anchor=west, xshift=1.0cm, omprofile, width=0.46\linewidth, height=4.6cm,
  xmin=0, xmax=10, ymin=0, ymax=10, xlabel={coordenada de reação}, ylabel={energia},
  title={\small A: $\mathrm{ML_5X} + \mathrm Y \to \mathrm{ML_5XY} \to \mathrm{ML_5Y}$}]
\addplot[omThm, very thick, smooth] coordinates {(0,2) (1.5,2.1) (3,6.5) (4.2,4.4) (5.0,4.2) (5.8,4.4) (7,7.2) (8.5,1.6) (10,1.5)};
\node[font=\scriptsize, anchor=north, align=center] at (axis cs:5,4.0) {ML$_5$XY (NC 7)};
\end{axis}
\end{tikzpicture}

{\small Perfis de energia de uma substituição dissociativa e de uma associativa
(esquemáticos): ambas passam por um intermediário, de número de coordenação (NC)
menor em D e de número de coordenação maior em A. Nos \omterm{def:b3:complex-mechanisms:mechanisms}{mecanismos de intercâmbio},
o poço desaparece e resta um único estado de transição.}
\end{omfigure}

Um complexo octaédrico que se substitui por uma pirâmide de base quadrada em geral
conserva sua configuração ($cis$ continua $cis$); um que passa por uma bipirâmide
trigonal, como muitos complexos de cobalto(III) na hidrólise básica, pode dar uma
mistura de produtos $cis$ e $trans$.

\section{Substituição em complexos quadrado-planares e efeito trans}

Os complexos quadrado-planares de íons $d^8$ (\ce{Pt^{2+}}, \ce{Pd^{2+}}, \ce{Au^{3+}}, \ce{Ni^{2+}} com ligantes fortes)
são coordenativamente insaturados: um ligante de entrada pode se aproximar pela
direção axial vazia. Suas substituições são associativas, por um estado de
transição pentacoordenado bipiramidal trigonal.

\begin{proposition}[Lei de velocidade de dois termos]\label{prop:b3:complex-mechanisms:square-planar}
A substituição $\mathrm{[PtL_3X] + Y \to [PtL_3Y] + X}$ em um solvente coordenante S segue
$k_{\mathrm{obs}} = k_1 + k_2[\mathrm Y]$ com Y em excesso.
\end{proposition}

\begin{proof}
Dois caminhos associativos paralelos: o ataque direto de Y ($k_2[\mathrm Y]$) e o ataque do
solvente, presente em grande excesso (pseudoprimeira ordem, $k_1$), que dá $[\mathrm{PtL_3S}]$, depois
rapidamente convertido por Y. Caminhos paralelos de primeira ordem se somam.
\end{proof}

\begin{definition}[Efeito trans, influência trans]\label{def:b3:complex-mechanisms:trans-effect}
O \emph{efeito trans}\index{efeito trans} é a aceleração da substituição
de um ligante pelo ligante em posição trans a ele: um efeito cinético, sobre o estado de transição.
A \emph{influência trans}\index{influencia trans@influência trans} é o enfraquecimento, no estado
fundamental, da ligação trans a um ligante, visto nos comprimentos de ligação e nas
frequências vibracionais: um efeito termodinâmico.
\end{definition}

O \omterm{def:b3:complex-mechanisms:trans-effect}{efeito trans} segue a ordem $\ce{CO}$, $\ce{CN-}$, $\ce{C2H4} > \ce{PR3}$, $\ce{H-} > \ce{CH3-} > \ce{I-}$, $\ce{SCN-} >
\ce{Br-} > \ce{Cl-} > \ce{NH3}$, $\ce{H2O}$. Os $\sigma$-doadores fortes enfraquecem a ligação trans (\omterm{def:b3:complex-mechanisms:trans-effect}{influência trans});
os $\pi$-aceptores fortes estabilizam o estado de transição pentacoordenado retirando
\omterm{def:b3:computational-chemistry:dft}{densidade eletrônica} do metal. Ambos aumentam a velocidade.

\begin{method}[Planejar a síntese de isômeros de platina(II)]\label{met:b3:complex-mechanisms:pt-synthesis}
\begin{enumerate}
  \item A cada etapa, o ligante substituído é o que está em trans ao ligante de
        maior \omterm{def:b3:complex-mechanisms:trans-effect}{efeito trans} presente.
  \item Ordene as adições de modo que essa regra leve ao isômero desejado.
\end{enumerate}
\end{method}

\begin{omfigure}
\setchemfig{atom sep=2.2em}
\schemestart
\chemfig{Cl-Pt(-[2]Cl)(-[6]Cl)-Cl}
\arrow{->[\ce{NH3}]}
\chemfig{Cl-Pt(-[2]Cl)(-[6]Cl)-NH_3}
\arrow{->[\ce{NH3}]}
\chemfig{Cl-Pt(-[2]NH_3)(-[6]Cl)-NH_3}
\schemestop

\bigskip
\schemestart
\chemfig{H_3N-Pt(-[2]NH_3)(-[6]NH_3)-NH_3}
\arrow{->[\ce{Cl-}]}
\chemfig{H_3N-Pt(-[2]NH_3)(-[6]NH_3)-Cl}
\arrow{->[\ce{Cl-}]}
\chemfig{Cl-Pt(-[2]NH_3)(-[6]NH_3)-Cl}
\schemestop

{\small O \omterm{def:b3:complex-mechanisms:trans-effect}{efeito trans} em ação (cargas omitidas). Acima: a partir de $\ce{[PtCl4]^{2-}}$, a segunda
amônia substitui um cloreto em trans a um cloreto (Cl$^-$ tem o \omterm{def:b3:complex-mechanisms:trans-effect}{efeito trans} maior),
dando o isômero $cis$, a cisplatina. Abaixo: a partir de $\ce{[Pt(NH3)4]^{2+}}$, o segundo cloreto
substitui a amônia em trans ao primeiro cloreto, dando o isômero $trans$.}
\end{omfigure}

\section{Transferência de elétrons}

\begin{definition}[Transferência de elétrons de esfera externa e de esfera interna]\label{def:b3:complex-mechanisms:electron-transfer}
Na \emph{transferência de elétrons de esfera externa}\index{transferencia de eletrons de esfera externa@transferência de elétrons de esfera externa}, o
elétron passa entre dois complexos cujas esferas de coordenação permanecem intactas;
na \emph{transferência de elétrons de esfera interna}\index{transferencia de eletrons de esfera interna@transferência de elétrons de esfera interna}, ele
passa por um ligante que faz ponte entre os dois metais em um complexo binuclear
transitório. Uma \emph{reação de autotroca}\index{reacao de autotroca@reação de autotroca} é uma
transferência de elétrons entre os dois estados de oxidação de um mesmo par, sem
mudança química líquida, como $\ce{[Fe(H2O)6]^{3+} + [Fe(H2O)6]^{2+}}$.
\end{definition}

Henry Taube demonstrou o caminho de esfera interna em 1953 com uma reação escolhida para
que o produto contasse sua história. O cobalto(III) é inerte, o cobalto(II) \omterm{def:b3:complex-mechanisms:labile}{lábil}; o crômio(II)
é \omterm{def:b3:complex-mechanisms:labile}{lábil}, o crômio(III) inerte. Quando $\ce{[Co(NH3)5Cl]^{2+}}$ é reduzido por $\ce{[Cr(H2O)6]^{2+}}$ em
meio ácido, todo o crômio(III) formado carrega o cloreto, como $\ce{[Cr(H2O)5Cl]^{2+}}$, mesmo em
presença de cloreto livre: o cloreto precisa ter feito ponte entre os dois metais,
ter ficado ligado ao crômio(III) inerte assim que o elétron passou e ter sido
liberado pelo cobalto(II) \omterm{def:b3:complex-mechanisms:labile}{lábil}.

\begin{omfigure}
\begin{tikzpicture}[font=\small]
  \begin{scope}[scale=0.8]
    \pic[transform shape] (a) at (0,0) {omoctahedron};
    \pic[transform shape] (b) at (2.0,0) {omoctahedron};
  \end{scope}
  \node[anchor=north east] at (0.3,-1.0) {$\mathrm{Co^{III}(NH_3)_5}$};
  \node[anchor=north west] at (1.3,-1.0) {$\mathrm{Cr^{II}(H_2O)_5}$};
  \node[omThm, anchor=south] at (0.8,0.15) {Cl};
  \draw[->, thick, omDef] (1.6,1.75) to[bend right=40] node[midway, above] {e$^-$} (0,1.75);
\end{tikzpicture}

{\small O intermediário em ponte de Taube (esquemático): o cloreto, ainda ligado ao
cobalto(III), entra na esfera de coordenação do crômio(II) substituindo uma molécula
de água; o elétron atravessa a ponte, e o cloreto fica no
crômio(III), agora inerte.}
\end{omfigure}

Um elétron salta em cerca de $10^{-16}$ s, muito mais depressa do que os núcleos se movem: pelo
\omterm{thm:b3:electronic-spectroscopy:franck-condon}{princípio de Franck--Condon} (\cref{ch:b3:electronic-spectroscopy}), a transferência
só pode ocorrer em uma configuração nuclear em que reagentes e produtos têm
a mesma energia. Na autotroca do ferro, as ligações \ce{Fe^{III}}--O são mais curtas que
as ligações \ce{Fe^{II}}--O; antes que o elétron possa passar, os dois complexos precisam se distorcer até
uma geometria intermediária comum, ao custo de energia.

\section{Teoria de Marcus}

\begin{definition}[Energia de reorganização]\label{def:b3:complex-mechanisms:reorganisation}
A \emph{energia de reorganização}\index{energia de reorganizacao@energia de reorganização} $\lambda$ de uma
transferência de elétrons é a energia necessária para levar os reagentes, sem
transferir o elétron, à configuração nuclear (comprimentos de ligação dos
complexos e orientações do solvente ao redor) dos produtos. Ela é a
soma de uma parte interna, das ligações, e de uma parte externa, do solvente.
\end{definition}

\begin{theorem}[Equação de Marcus]\label{thm:b3:complex-mechanisms:marcus}
Se as energias de Gibbs dos reagentes e dos produtos são parábolas de mesma
curvatura em uma coordenada de reação $x$ (0 no equilíbrio dos reagentes, 1 no dos
produtos), a \omterm{def:b3:rate-theories:activation}{energia de Gibbs de ativação} de uma transferência de elétrons com
energia de Gibbs padrão de reação $\Delta G^\circ$ e \omterm{def:b3:complex-mechanisms:reorganisation}{energia de reorganização} $\lambda$ é
\[
  \Delta G^{\ddagger} = \frac{(\lambda + \Delta G^\circ)^2}{4\lambda} .
\]
Esta é a \emph{equação de Marcus}\index{equacao de Marcus@equação de Marcus}.
\end{theorem}

\begin{proof}
Escreva $G_R = \lambda x^2$ e $G_P = \lambda(x - 1)^2 + \Delta G^\circ$: a curvatura é fixada por $G_R(1) = \lambda$, a
definição de $\lambda$. A transferência ocorre onde as curvas se cruzam: $\lambda x^2 = \lambda x^2 - 2\lambda x + \lambda +
\Delta G^\circ$, logo $x^\ast = (\lambda + \Delta G^\circ)/2\lambda$, e $\Delta G^{\ddagger} = G_R(x^\ast) = (\lambda + \Delta G^\circ)^2/4\lambda$.
\end{proof}

\begin{corollary}[Região invertida]\label{cor:b3:complex-mechanisms:inverted}
Com $\lambda$ fixo, a velocidade aumenta à medida que a reação se torna mais exergônica até
$-\Delta G^\circ = \lambda$, onde $\Delta G^{\ddagger} = 0$; além disso, ela volta a diminuir.
\end{corollary}

\begin{proof}
$\Delta G^{\ddagger}$ é uma parábola em $\Delta G^\circ$ com mínimo, zero, em $\Delta G^\circ = -\lambda$; para $\Delta G^\circ < -\lambda$, ela volta a
crescer.
\end{proof}

\begin{definition}[Região invertida]\label{def:b3:complex-mechanisms:inverted}
A \emph{região invertida}\index{regiao invertida@região invertida} da transferência de elétrons é a
faixa de forças motrizes $-\Delta G^\circ > \lambda$ na qual a velocidade diminui à medida que a reação se torna
mais exergônica.
\end{definition}

\begin{omfigure}
\begin{tikzpicture}
\begin{axis}[name=mp, width=0.47\linewidth, height=5.6cm, xmin=-0.6, xmax=2.2, ymin=-160, ymax=120,
  axis lines=left, xlabel={coordenada de reação $x$}, ylabel={$G$ / \unit{kJ.mol^{-1}}},
  tick label style={font=\small}, label style={font=\small}, xtick={0,1,2}]
\addplot[black, thick] table[x=x, y=GR] {figdata/out/bachelor-3-marcus-parabolas.dat};
\addplot[omDef!50, thick] table[x=x, y=P0] {figdata/out/bachelor-3-marcus-parabolas.dat};
\addplot[omDef, thick] table[x=x, y=P50] {figdata/out/bachelor-3-marcus-parabolas.dat};
\addplot[omThm, thick] table[x=x, y=P100] {figdata/out/bachelor-3-marcus-parabolas.dat};
\addplot[omThm!60, thick, dashed] table[x=x, y=P150] {figdata/out/bachelor-3-marcus-parabolas.dat};
\node[font=\scriptsize, anchor=north] at (axis cs:0,-2) {R};
\end{axis}
\begin{axis}[at={(mp.east)}, anchor=west, xshift=1.4cm, width=0.47\linewidth, height=5.6cm,
  xmin=0, xmax=250, ymin=0, ymax=11.5, axis lines=left, xlabel={$-\Delta G^\circ$ / \unit{kJ.mol^{-1}}},
  ylabel={$\log(k/\unit{s^{-1}})$}, tick label style={font=\small}, label style={font=\small}]
\addplot[omDef, very thick] table[x=mdG, y=logk] {figdata/out/bachelor-3-marcus-rate.dat};
\draw[black!40, dashed] (axis cs:100,0) -- (axis cs:100,11);
\node[font=\scriptsize, anchor=south east] at (axis cs:95,1) {normal};
\node[font=\scriptsize, anchor=south west] at (axis cs:105,1) {invertida};
\end{axis}
\end{tikzpicture}

{\small Teoria de Marcus com $\lambda = \qty{100}{kJ/mol}$ (modelo). À esquerda: a parábola dos reagentes
(preto) e as parábolas dos produtos para $\Delta G^\circ = 0$ e $-50$ (azul), $-100 = -\lambda$ (vermelho) e
\qty{-150}{kJ/mol} (tracejado, \omterm{def:b3:complex-mechanisms:inverted}{região invertida}); seu ponto de
cruzamento, o estado de transição, desce até o fundo da curva dos reagentes em
$-\Delta G^\circ = \lambda$ e volta a subir além disso. À direita: o logaritmo da constante de velocidade
(fator pré-exponencial do modelo \qty{e11}{s^{-1}}) passa por um máximo em $-\Delta G^\circ = \lambda$.}
\end{omfigure}

\begin{proposition}[Energia de reorganização de esfera externa]\label{prop:b3:complex-mechanisms:outer-reorganisation}
Para dois reagentes esféricos de raios $a_1$, $a_2$ a uma distância $d$ em um solvente de
índice de refração $n$ e permissividade relativa $\varepsilon_s$, a parte de $\lambda$ devida ao solvente é
proporcional a $\bigl(\frac{1}{2a_1} + \frac{1}{2a_2} - \frac1d\bigr)\bigl(\frac{1}{n^2} - \frac{1}{\varepsilon_s}\bigr)$.
\end{proposition}

\admitted

A água, muito polar ($\varepsilon_s$ grande) mas com um índice de refração comum, dá uma grande
reorganização do solvente; os solventes apolares, uma pequena. Reagentes grandes
se reorganizam menos: é por isso que as proteínas de transferência de elétrons enterram seus sítios metálicos
no interior da proteína.

\begin{theorem}[Relação cruzada de Marcus]\label{thm:b3:complex-mechanisms:cross-relation}
Para uma reação cruzada $\mathrm{A_{ox} + B_{red} \to A_{red} + B_{ox}}$ de constante de equilíbrio $K_{12}$, entre pares
cujas constantes de velocidade de autotroca são $k_{11}$ e $k_{22}$,
\[
  k_{12} = \sqrt{k_{11}k_{22}K_{12}f}, \qquad \ln f = \frac{(\ln K_{12})^2}{4\ln(k_{11}k_{22}/Z^2)},
\]
com $Z$ o fator de frequência de colisão; $f \approx 1$ quando $K_{12}$ não é grande demais. Esta é a
\emph{relação cruzada de Marcus}\index{relacao cruzada de Marcus@relação cruzada de Marcus}.
\end{theorem}

\begin{proof}[Demonstração parcial]
Suponha que a \omterm{def:b3:complex-mechanisms:reorganisation}{energia de reorganização} da reação cruzada seja a média das das
autotrocas, $\lambda_{12} = (\lambda_{11} + \lambda_{22})/2$ (cada reagente contribui com metade da reorganização de sua
própria autotroca), e que todas as constantes de velocidade sejam $k = Z\eu^{-\Delta G^{\ddagger}/RT}$. Então
$\Delta G_{12}^{\ddagger} = \frac{\lambda_{12}}{4}\bigl(1 + \frac{\Delta G^\circ_{12}}{\lambda_{12}}\bigr)^2 = \frac{\lambda_{12}}{4} + \frac{\Delta G^\circ_{12}}{2} + \frac{(\Delta G^\circ_{12})^2}{4\lambda_{12}}$. O primeiro termo dá
$\sqrt{k_{11}k_{22}}$ (pois $\Delta G^{\ddagger}_{ii} = \lambda_{ii}/4$), o segundo $\sqrt{K_{12}}$ (com $\Delta G^\circ_{12} = -RT\ln K_{12}$), o terceiro
o fator $f$.
\end{proof}

\begin{method}[Velocidade de uma reação cruzada a partir de dados de autotroca]\label{met:b3:complex-mechanisms:marcus-estimate}
\begin{enumerate}
  \item Encontre as constantes de autotroca $k_{11}$, $k_{22}$ dos dois pares.
  \item Calcule $K_{12}$ a partir dos potenciais padrão: $\ln K_{12} = nF(E^\circ_{\text{oxidante}} - E^\circ_{\text{redutor}})/RT$.
  \item $k_{12} = \sqrt{k_{11}k_{22}K_{12}}$ (com $f$ se $K_{12}$ for muito grande). Uma concordância dentro de um fator
        de poucas unidades com o valor medido indica um mecanismo de esfera externa.
\end{enumerate}
\end{method}

Marcus publicou a teoria em 1956; a \omterm{def:b3:complex-mechanisms:inverted}{região invertida}, de início posta em dúvida, foi
observada em 1984 por Gerhard Closs e John Miller em moléculas nas quais doador e
aceptor são mantidos a uma distância fixa por um espaçador rígido, de modo que a difusão
não pode escondê-la.

\begin{inthelab}[Troca de água por RMN de oxigênio-17]
A ressonância de oxigênio-17 da água ligada a um íon paramagnético e a da água
livre são alargadas pela troca entre elas. Registrar a largura de linha do
sinal da água livre em função da temperatura, em uma solução do perclorato do metal
enriquecida em \ce{^{17}O}, dá a constante de velocidade de troca e, em uma sonda de alta pressão,
seu \omterm{def:b3:complex-mechanisms:activation-volume}{volume de ativação}.
\end{inthelab}

\begin{safety}{\ghs[0.8cm]{GHS05}\ghs[0.8cm]{GHS06}\ghs[0.8cm]{GHS08}}
% ledger: ghs4:K2PtCl4
Tetracloroplatinato(II) de potássio: tóxico por ingestão, provoca lesões oculares
graves, sensibilizante respiratório e cutâneo. Os complexos de platina são manipulados com
luvas em capela; a cisplatina e seus análogos são fármacos citotóxicos e são
manipulados como tais.
\end{safety}

\begin{history}[Taube e Marcus]
Henry Taube classificou os complexos em \omterm{def:b3:complex-mechanisms:labile}{lábeis} e inertes em 1952 e, com sua
experiência crômio--cobalto de 1953, estabeleceu o mecanismo de esfera interna;
recebeu o Prêmio Nobel de Química de 1983. Rudolph Marcus deduziu, a partir de 1956, a
teoria da transferência de elétrons que leva seu nome, e recebeu o Prêmio Nobel de Química
de 1992.
\end{history}

\section{Exercícios}

\begin{exercise}[$\star$]\label{exo:b3:complex-mechanisms:1}
Classifique como \omterm{def:b3:complex-mechanisms:labile}{lábil} ou inerte, justificando: $\ce{[Cr(H2O)6]^{3+}}$, $\ce{[Co(NH3)6]^{3+}}$, $\ce{[Cr(H2O)6]^{2+}}$,
$\ce{[Cu(H2O)6]^{2+}}$, $\ce{[Zn(H2O)6]^{2+}}$.
\end{exercise}

\begin{exercise}[$\star$]\label{exo:b3:complex-mechanisms:2}
Mostre que a lei de velocidade D se reduz a $v = k_1[\mathrm{ML_5X}]$ para $[\mathrm Y]$ alto e a $v = (k_1k_2/k_{-1})
[\mathrm{ML_5X}][\mathrm Y]/[\mathrm X]$ quando $k_{-1}[\mathrm X] \gg k_2[\mathrm Y]$.
\end{exercise}

\begin{exercise}[$\star$]\label{exo:b3:complex-mechanisms:3}
Preveja os sinais de $\Delta V^{\ddagger}$ e de $\Delta^{\ddagger}S^\circ$ para uma substituição D e uma A.
\end{exercise}

\begin{exercise}[$\star$]\label{exo:b3:complex-mechanisms:4}
Dê os produtos de $\ce{[PtCl4]^{2-}}$ com dois \ce{NH3} e de $\ce{[Pt(NH3)4]^{2+}}$ com dois \ce{Cl-}.
\end{exercise}

\begin{exercise}[$\star\star$]\label{exo:b3:complex-mechanisms:5}
Para uma substituição em um aquaíon, $K_{\mathrm{os}} = \qty{2.0}{L.mol^{-1}}$ e $k_i = \qty{3.0e4}{s^{-1}}$ (dados do
exercício). Calcule $k_{\mathrm{obs}}$ para $[\mathrm Y] = 0.10$ e \qty{1.0}{mol/L}, e seu limite.
\end{exercise}

\begin{exercise}[$\star\star$]\label{exo:b3:complex-mechanisms:6}
Uma constante de velocidade dobra quando a pressão sobe de 0.1 para \qty{100}{MPa} a \qty{298}{K}.
Calcule $\Delta V^{\ddagger}$ e conclua.
\end{exercise}

\begin{exercise}[$\star\star$]\label{exo:b3:complex-mechanisms:7}
No $\ce{[PtCl3(PEt3)]-}$, a ligação Pt--Cl em trans ao \ce{PEt3} mede \qty{2.42}{\angstrom} e as outras duas \qty{2.32}{\angstrom}
(dados do exercício). Interprete e preveja qual cloreto é substituído primeiro.
\end{exercise}

\begin{exercise}[$\star\star$]\label{exo:b3:complex-mechanisms:8}
Liste as evidências experimentais que mostrariam que uma transferência de elétrons é de
esfera interna.
\end{exercise}

\begin{exercise}[$\star\star$]\label{exo:b3:complex-mechanisms:9}
Para $\lambda = \qty{80}{kJ/mol}$, calcule $\Delta G^{\ddagger}$ para $\Delta G^\circ = 0$, $-20$, $-80$ e \qty{-120}{kJ/mol}.
\end{exercise}

\begin{exercise}[$\star\star\star$]\label{exo:b3:complex-mechanisms:10}
Dois pares têm constantes de autotroca $k_{11} = 4.0$ e $k_{22} = \qty{2.0e5}{L.mol^{-1}.s^{-1}}$, e
sua reação cruzada tem $K_{12} = \num{1.0e4}$ (dados do exercício). Estime $k_{12}$ ($f = 1$).
\end{exercise}

\begin{exercise}[$\star\star\star$]\label{exo:b3:complex-mechanisms:11}
Usando o modelo só $\sigma$, calcule as energias de ativação do campo ligante dos
íons $d^5$, $d^7$ e $d^8$ de spin alto e ordene \ce{Mn^{2+}}, \ce{Co^{2+}} e \ce{Ni^{2+}} pela
velocidade de troca de água esperada.
\end{exercise}

\begin{exercise}[$\star\star\star$]\label{exo:b3:complex-mechanisms:12}
Por que a \omterm{def:b3:complex-mechanisms:inverted}{região invertida} é difícil de observar em reações entre íons que se
encontram por difusão em solução, e como as moléculas rígidas doador--aceptor a
revelaram?
\end{exercise}

\section{Problema: Fazer cisplatina, e não transplatina}

\begin{problem}[{Problema de fim de semana --- fazer cisplatina, e não transplatina: o
íon platina(II) quadrado-planar, uma síntese planejada com o efeito trans, a
aquação do fármaco e por que ele é ativado dentro das células}]
\label{pb:b3:complex-mechanisms:1}
% ledger: ghs4:K2PtCl4
A cisplatina, $cis$-$\ce{[PtCl2(NH3)2]}$, é preparada a partir de $\ce{K2[PtCl4]}$. Dados do problema: a
\qty{37}{\celsius}, a primeira aquação, $\ce{[PtCl2(NH3)2] + H2O -> [PtCl(H2O)(NH3)2]+ + Cl-}$, tem
constante de velocidade $k = \qty{9.0e-5}{s^{-1}}$ e constante de equilíbrio $K = \qty{3.0e-3}{mol/L}$; a
concentração de cloreto é de cerca de \qty{100}{mM} no plasma sanguíneo e de \qty{4}{mM} dentro das células.

\medskip\noindent\textbf{Parte I --- A platina(II).}
\begin{enumerate}
  \item Dê o número de elétrons $d$ de \ce{Pt^{2+}}.
  \item Por que seus complexos são quadrado-planares e diamagnéticos?
  \item Por que suas substituições são associativas?
  \item Escreva a lei de velocidade de dois termos e explique seus dois termos.
  \item O fármaco é \omterm{def:b3:complex-mechanisms:labile}{lábil} ou inerte na escala de tempo de um dia?
\end{enumerate}

\medskip\noindent\textbf{Parte II --- A síntese.}
\begin{enumerate}[resume]
  \item O que $\ce{K2[PtCl4]}$ dá com dois \ce{NH3}? Por quê?
  \item Essa rota direta dá um produto impuro. A rota de Dhara converte primeiro
        $\ce{[PtCl4]^{2-}}$ em $\ce{[PtI4]^{2-}}$ com excesso de KI. Por que o iodeto é útil?
  \item A adição de \ce{NH3} dá $cis$-$\ce{[PtI2(NH3)2]}$. Explique a estereoquímica.
  \item Os iodetos são então removidos com nitrato de prata em água. O que se forma?
  \item A adição de KCl dá a cisplatina. Por que a configuração não muda?
  \item Como você prepararia a transplatina?
  \item Por que a transplatina não serve como fármaco, embora também reaja?
\end{enumerate}

\medskip\noindent\textbf{Parte III --- A aquação.}
\begin{enumerate}[resume]
  \item Calcule a meia-vida da primeira aquação.
  \item Calcule a fração aquada após \qty{2.0}{h} em uma solução sem cloreto.
  \item Por que o aquacomplexo é a forma reativa em relação ao DNA?
  \item A aquação é dissociativa ou associativa? Qual seria seu \omterm{def:b3:complex-mechanisms:activation-volume}{volume de
        ativação}?
  \item Por que uma solução injetável do fármaco é preparada em soro fisiológico?
\end{enumerate}

\medskip\noindent\textbf{Parte IV --- Plasma e célula.}
\begin{enumerate}[resume]
  \item Escreva a fração de equilíbrio do aquacomplexo em função de
        $[\ce{Cl-}]$.
  \item Calcule-a no plasma.
  \item Calcule-a dentro de uma célula.
  \item Calcule a razão.
  \item Por que, então, o fármaco age sobretudo dentro das células?
  \item Que outros ligantes dentro de uma célula podem se ligar à platina e desativar o
        fármaco?
  \item Por que os dois ligantes amin permanecem ligados o tempo todo?
  \item Enuncie o resultado: a razão entre a fração de aquacomplexo dentro de uma célula e
        a fração no plasma.
\end{enumerate}
\end{problem}
